\subsection{\texorpdfstring{Differentials and \(p\)-bases}{Differentials and p-bases}}

Let K be a field (of characteristic p) and V a vector space over K. We recall (III, p.~552) that a derivation of K into V is a mapping D of K into V satisfying the relations

\begin{enumerate}
\def\labelenumi{(\arabic{enumi})}
\tightlist
\item
\end{enumerate}

\[
\begin{array}{r l} \mathrm{D} (x + y) & = \mathrm{D} (x) + \mathrm{D} (y) \\ \mathrm{D} (x y) & = x \mathrm{D} (y) + y \mathrm{D} (x) \end{array}\tag{2}
\]

for \(x, y \in \mathbf{K}\) . It follows that \(\mathrm{D}(1) = 0\) and \(\mathrm{D}(x^n) = nx^{n-1}\mathrm{D}(x)\) for all \(x \neq 0\) and all \(n \in \mathbf{Z}\) (III, p.~557 and 558). Since \(\mathbf{K}\) is of characteristic \(p\) , we thus have for all \(x \in \mathbf{K}\)

\[
\mathrm{D} \left(x ^ {p}\right) = 0.\tag{3}
\]

Moreover (III, p.~558), we have for \(x\) , \(y \in \mathbf{K}\) , \(y \neq 0\) ,

\[
\mathrm{D} (x / y) = \left(y \mathrm{D} (x) - x \mathrm{D} (y)\right) / y ^ {2}.\tag{4}
\]

By the above formulae the kernel of D is a subfield E of K containing \(K^{p}\) . Let M be a subfield of K; we say that D is an M-derivation if it is M-linear; by (2) it comes to the same to assume that the restriction of D to M is zero.

We have (III, p.~570) defined the module \(\Omega_{\mathrm{M}}(\mathbf{K})\) of M-differentials of K and the canonical M-derivation \(d = d_{\mathrm{K / M}}\) of K into \(\Omega_{\mathrm{M}}(\mathbf{K})\) . The image of \(d_{\mathrm{K / M}}\) generates the vector K-space \(\Omega_{\mathrm{M}}(\mathbf{K})\) . For a mapping D of K into V to be an M-derivation it is necessary and sufficient that there should exist a K-linear mapping \(u: \Omega_{\mathrm{M}}(\mathbf{K}) \to \mathbf{V}\) such that \(\mathbf{D} = u \circ d_{\mathrm{K / M}}\) ; this mapping \(u\) is uniquely determined.

PROPOSITION 5. --- Let K be a field and L an extension of K generated by an element x such that \(x \notin K\) , \(x^{p} \in K\) . Let V be a vector space over L and \(\Delta\) a derivation of K into V such that \(\Delta(x^{p}) = 0\) . Then there exists a unique derivation D of L into V extending \(\Delta\) and such that D(x) = 0.

By V, p.~24, Prop. 1, \(\{1, x, \ldots, x^{p-1}\}\) is a basis of L over K. For every element \(u = a_{0} + a_{1}x + \cdots + a_{p-1}x^{p-1}\) of L (with \(a_{0}, \ldots, a_{p-1}\) in K) we put \(\mathrm{D}(u) = \Delta(a_{0}) + x\Delta(a_{1}) + \cdots + x^{p-1}\Delta(a_{p-1})\) . It is clear that D extends \(\Delta\) and satisfies (1); so it is enough to establish the relation

\[
\mathrm{D} (u v) = u \cdot \mathrm{D} (v) + v \cdot \mathrm{D} (u)
\]

when \(u\) is of the form \(ax^i\) and \(v\) of the form \(bx^j\) with \(a, b\) in \(\mathbf{K}\) , \(0 \leqslant i < p\) and \(0 \leqslant j < p\) .

When \(i + j < p\) we have \(uv = x^{i + j} \cdot ab\) , whence \(D(uv) = x^{i + j}\Delta(ab)\) . Otherwise we have \(0 \leqslant i + j - p < p\) and so \(uv = x^{i + j - p}(abx^p)\) , with \(abx^p \in K\) , but since \(\Delta(x^p) = 0\) , we have \(\Delta(abx^p) = x^p\Delta(ab)\) , whence again \(D(uv) = x^{i + j}\Delta(ab)\) . So we have in all cases

\[
\begin{array}{r l} \mathrm{D} (u v) & = x ^ {i + j} \Delta (a b) = x ^ {i} x ^ {j} (a \cdot \Delta (b) + b \cdot \Delta (a)) \\ & = (a x ^ {i}) x ^ {j} \cdot \Delta (b) + (b x ^ {j}) \cdot x ^ {i} \cdot \Delta (a) = u \cdot \mathrm{D} (v) + v \cdot \mathrm{D} (u). \end{array}
\]

COROLLARY 1. --- Let K be a field, L a p-radical extension of height ≤ 1 of K and V a vector space over L. Every derivation Δ of K into V vanishing on L \(^{p}\) extends to a derivation of L into V.

By Zorn's lemma (E, III, p.~20) there exists a maximal extension of \(\Delta\) to a derivation \(D_0: L_0 \to V\) , where \(L_0\) is a subfield of \(L\) containing \(K\) . Let \(x \in L\) ; we have \(x^p \in L_0\) and \(\Delta(x^p) = 0\) , whence \(D_0(x^p) = 0\) . By Prop. 5,

\(D_{0}\) extends to a derivation defined on \(L_{0}(x)\) ; by the maximal character of \(D_{0}\) we thus have \(L_{0}(x) = L_{0}\) whence \(x \in L_{0}\) . Since x was arbitrary we conclude that \(L_{0} = L\) .

COROLLARY 2. --- Let L be a p-radical extension of height ≤ 1 of a field K and let E be a subextension of L. Let U be the subspace of \(\Omega_{\mathrm{K}}(\mathrm{L})\) generated by the differentials of elements of E. Then E consists of the elements of L whose differential belongs to U. In particular we have \(d_{L/K}x \neq 0\) for all \(x \in L - K\) .

Let \(x\) be an element of \(L\) not belonging to \(E\) . Then \(\{1, x, \ldots, x^{p-1}\}\) is a basis of \(E(x)\) over \(E\) . Let \(\Delta\) be the \(E\) -linear mapping of \(E(x)\) into \(L\) such that \(\Delta(x^i) = ix^i\) for \(0 \leqslant i < p\) . We verify at once that \(\Delta\) is an \(E\) -derivation of \(E(x)\) into \(L\) , and in particular a \(K\) -derivation. By Cor. 1 to Prop. 5 there exists a \(K\) -derivation \(D\) of \(L\) into \(L\) extending \(\Delta\) . By the universal property of \(\Omega_K(L)\) , there exists a linear form \(u\) on this vector \(L\) -space such that \(D = u \circ d_{L/K}\) . We have \(D(x) = x\) and \(D|E = 0\) ; therefore we have \(u(d_{L/K}x) \neq 0\) and \(u|U = 0\) , whence \(d_{L/K}x \notin U\) .

The last assertion is the particular case E = K; then we have U = 0.

THEOREM 1. --- Let L be a p-radical extension of height ≤ 1 of a field K and \((a_{i})_{i \in I}\) a family of elements of L.

\begin{enumerate}
\def\labelenumi{\alph{enumi})}
\item
  For \((x_{i})_{i\in I}\) to be \(p\) -free over \(\mathbf{K}\) it is necessary and sufficient that the family \((dx_{i})_{i\in I}\) should be free in the vector \(\mathbf{L}\) -space \(\Omega_{\mathbf{K}}(\mathbf{L})\) .
\item
  For \((x_{i})_{i\in I}\) to generate L over K it is necessary and sufficient that the family \((dx_{i})_{i\in I}\) should generate the vector L-space \(\Omega_{\mathbf{K}}(\mathbf{L})\) .
\item
  For \((x_{i})_{i\in I}\) to be a \(p\) -basis of \(\mathbf{L}\) over \(\mathbf{K}\) it is necessary and sufficient that the family \((dx_{i})_{i\in I}\) should be a basis of the vector \(\mathbf{L}\) -space \(\Omega_{\mathbf{K}}(\mathbf{L})\) .
\end{enumerate}

To begin with let us remark that the differential \(dx\) of an element of the field \(\mathrm{K}((x_i)_{i \in I})\) is a linear combination with coefficients in \(L\) of the differentials \(dx_i\) , \(i \in I\) . For the family \((x_i)_{i \in I}\) to be \(p\) -free it is necessary and sufficient that \(x_i \notin \mathrm{K}(x_j)_{j \in I - \{i\}}\) for all \(i \in I\) (V, p.~99, Prop. 2). By Cor. 2 of Prop. 5 this means that \(dx_i\) is not a linear combination of the \(dx_j\) for \(j \neq i\) in \(I\) . Assertion \(a)\) follows from this. Now \(b)\) follows at once from Cor. 2 to Prop. 5 and \(c)\) follows from \(a)\) and \(b)\) .

The following corollary makes Cor. 1 of Prop. 5 more specific :

COROLLARY. --- Let \((x_{i})_{i\in I}\) be a p-basis of L over K. Let V be a vector space over L, \(\Delta\) a derivation of K into V which vanishes on \(L^{p}\) and \((u_{i})_{i\in I}\) a family of elements of V. Then there exists a unique derivation D of L in V which extends \(\Delta\) and maps \(x_{i}\) to \(u_{i}\) for all \(i\in I\) .

By Cor. 1 of Prop. 5 there exists a derivation \(D_{0}\) of L into V extending \(\Delta\) . The derivations of L in V extending \(\Delta\) are precisely the mappings of the form \(D = D_{0} + u \circ d_{L/K}\) where u is an L-linear mapping of \(\Omega_{\mathrm{K}}(\mathrm{L})\) into V. We have \(\mathrm{D}(x_{i}) = u_{i}\) if and only if the linear mapping u satisfies the conditions \(u(dx_{i}) = u_{i} - \mathrm{D}_{0}(x_{i})\) . Since the family \(dx_{i}\) is a basis of \(\Omega_{\mathrm{K}}(\mathrm{L})\) , this determines u uniquely and the corollary follows.

Let L be a p-radical extension of height \(\leqslant1\) of K and \((x_{i})_{i\in I}\) a p-basis of L over K. By the Cor. to Th. 1 there exists for each \(i\in I\) a K-derivation \(D_{i}\) of L into L characterized by \(\mathrm{D}_{i}(x_{j})=\delta_{ij}\) (Kronecker symbol); \(D_{i}\) is sometimes called the partial derivation with respect to \(x_{i}\) . When I is finite, the family \((\mathrm{D}_{i})_{i\in I}\) is a basis of the vector space over L consisting of the K-derivations of L into L.

Th. 1 allows the study of p-bases to be reduced to the study of bases of a vector space. For example we have the following result:

THEOREM 2. --- Let L be a p-radical extension of height ≤1 of a field K.

\begin{enumerate}
\def\labelenumi{\alph{enumi})}
\item
  There exist p-bases of L over K. More precisely, if S is a p-free subset of L over K and T a subset of L such that \(S \subset T\) and \(L = K(T)\) , then there exists a p-basis B of L over K such that \(S \subset B \subset T\) .
\item
  Two \(p\) -bases of \(\mathbf{L}\) over \(\mathbf{K}\) have the same cardinal.
\item
  For \([L:K]\) to be finite it is necessary and sufficient that the vector space \(\Omega_{K}(L)\) over L should be of finite dimension over L and then we have
\end{enumerate}

\[
[ \mathrm{L}: \mathrm{K} ] = p ^ {[ \Omega_ {\mathrm{K}} (\mathrm{L}): \mathrm{L} ]}.\tag{5}
\]

Assertions a) and b) are immediate (II, p.~292). If {[}L: K{]} is finite, there exists by a) a finite \(p\) -basis \((x_1, \ldots, x_n)\) of L over K; then the monomials \(x_1^{\alpha_1} \ldots x_n^{\alpha_n}\) with \(0 \leqslant \alpha_i < p\) for \(1 \leqslant i \leqslant n\) form a basis of L over K and the differentials \(dx_1, \ldots, dx_n\) form a basis of \(\Omega_K(L)\) over L. We thus have \([L:K] = p^n\) and \([\Omega_K(L):L] = n\) . Conversely if \(\Omega_K(L)\) is of finite dimension over L, there exists a finite \(p\) -basis of L over K (V, p.~102, Th. 1) and {[}L:K{]} is finite.

COROLLARY. --- For each \(x \in L - K\) , there exists a p-basis of L over K containing x.

For \(\{x\}\) is a \(p\) -free subset of \(\mathbf{L}\) over \(\mathbf{K}\) and so it is enough to apply Th. 2, \(a\) ). Let \(\mathbf{K}\) be a field and \(\mathbf{L}\) an extension of \(\mathbf{K}\) . If \(\mathbf{D}\) is a \(\mathbf{K}\) -derivation of \(\mathbf{L}\) with values in a vector \(\mathbf{L}\) -space, then we have \(\mathrm{D}(x^p) = 0\) for all \(x \in \mathbf{L}\) , and so \(\mathbf{D}\) is a \(\mathrm{K}(\mathrm{L}^p)\) -derivation; it follows that \(\Omega_{\mathrm{K}}(\mathrm{L}) = \Omega_{\mathrm{K}(\mathrm{L}^p)}(\mathrm{L})\) . Since \(\mathbf{L}\) is a \(p\) -radical extension of height \(\leqslant 1\) of \(\mathrm{K}(\mathrm{L}^p)\) , we can apply the above results. For example, from Th. 1, \(c\) ) we deduce the following: let \((x_i)_{i \in I}\) be a family of elements of \(\mathbf{L}\) ; for \((dx_i)_{i \in I}\) to be a basis of the vector \(\mathbf{L}\) -space \(\Omega_{\mathrm{K}}(\mathrm{L})\) it is necessary and sufficient that \((x_i)_{i \in I}\) should be a \(p\) -basis of \(\mathbf{L}\) over \(\mathrm{K}(\mathrm{L}^p)\) . Similarly, Cor. 2 of Prop. 5 implies:

PROPOSITION 6. --- Let K be a field (of characteristic p) and L an extension of K. a) Let \(x \in L\) ; the differential \(d_{L/K}x\) vanishes if and only if x belongs to \(\mathrm{K}(\mathrm{L}^{p})\) .

\begin{enumerate}
\def\labelenumi{\alph{enumi})}
\setcounter{enumi}{1}
\item
  \(\Omega_{\mathbf{K}}(\mathbf{L}) = 0\) if and only if \(\mathbf{L} = \mathbf{K}(\mathbf{L}^p)\) .
\item
  Suppose that \(\mathbf{L}\) is \(p\) -radical of finite height over \(\mathbf{K}\) ; then \(\Omega_{\mathbf{K}}(\mathbf{L}) = 0\) if and only if \(\mathbf{L} = \mathbf{K}\) .
\end{enumerate}

